\documentclass[dvipdfmx, uplatex]{jsarticle}
\usepackage[dvipdfmx]{graphicx}
\usepackage{amsmath}
\usepackage{amssymb}
\usepackage{amsfonts}
\usepackage{amsthm}
\usepackage{mathrsfs}
\usepackage{bm}
\usepackage{braket}

\usepackage{silence}
\WarningFilter*{hyperref}{Token not allowed in a PDF string}
\usepackage[dvipdfmx]{hyperref}
\usepackage{pxjahyper}
\allowdisplaybreaks[4]

\usepackage[left=25truemm,top=20truemm,bottom=20truemm,right=25truemm]{geometry}

\theoremstyle{definition}
\newtheorem{prob}{問題}

\title{演習：余剰次元（周期境界）上の拡散と次元簡約\\
\large ―― Kaluza--Klein 塔・ゼロモード・運動量--巻きつき双対 ――}
\author{}
\date{}

\begin{document}
\maketitle
\vspace{-6zh}

\begin{abstract}
前の演習（\texttt{exercise\_lattice\_random\_walk}）で，格子ランダムウォークの連続極限として
2 次元拡散方程式 $\partial_t P=D\nabla^2P$ を得た。本問では，$y$ 方向を半径 $R$ の円に
コンパクト化（周期境界条件）し，$y$ の Fourier モード $m$ ごとに分解する。
\textbf{ゼロモード}は $x$ 方向の純 1 次元拡散に帰着し（長時間・低エネルギーでの有効理論），
\textbf{非ゼロモード}は減衰率（``質量''）$\propto m^2/R^2$ を持つ Kaluza--Klein 的な塔を成して
短時間・高エネルギーで余剰次元の効果を担う。さらに，円上の熱核が\emph{運動量モード和}と
\emph{巻きつき（像）和}の 2 通りに表され（Poisson 和），互いに逆の極限で支配的になる――
拡散における運動量--巻きつき双対の構造を見る。
前演習（\texttt{exercise\_lattice\_random\_walk}）では連続極限で回転対称性 $O(2)$ が創発する
（異方的スケーリング $z=2$）ことを見たが，本問ではその\textbf{スケーリング次元}を計算する：
$x\to bx$ に対し $t\to b^2t$ と縮む異方的スケーリングのもとで，確率密度の $t$ 冪 $P\sim t^{-d/2}$ が
空間次元 $d$ を``測る''こと，そして本系では\emph{多数の KK モードが効く短時間で $t^{-1}$（2 次元的），
最低 KK モードだけが残る長時間で $t^{-1/2}$（1 次元的）}となって\textbf{有効次元が $2$（UV）から
$1$（IR）へ流れる}ことを見る。\textbf{各問は表式・数値を答えさせる形にしてあり，
本文中には原則として解答を与えていない}。完全な解答は別ファイル
\texttt{exercise\_diffusion\_extra\_dimension\_solutions.tex} にある。
\end{abstract}

\tableofcontents

%======================================================================
\section{設定}
%======================================================================
2 次元拡散方程式 $\partial_t P(x,y,t)=D\bigl(\partial_x^2+\partial_y^2\bigr)P$ を考える。
$x\in\mathbb{R}$ は無限に広いが，$y$ は\textbf{半径 $R$ の円}（周期 $L=2\pi R$，$y\sim y+2\pi R$）に
コンパクト化されているとする。初期条件は原点の点源 $P(x,y,0)=\delta(x)\,\delta(y)$
（$\delta(y)$ は円上の周期的 delta）。拡散係数 $D$ は前演習で得たもの。

%======================================================================
\section{問題}
%======================================================================

\begin{prob}\label{prob:kk}
\begin{enumerate}
\item \textbf{$y$ の Fourier モード分解。}周期性から $y$ 方向の波数は離散的になる。
許される波数 $k_y$ を整数 $m$ で表せ。$P(x,y,t)=\dfrac{1}{2\pi R}\displaystyle\sum_{m\in\mathbb{Z}}
P_m(x,t)\,e^{ik_y y}$ と展開したとき，各モード係数 $P_m(x,t)$ が満たす\textbf{1 次元の}偏微分方程式を導け。
その方程式の $x$ 方向の拡散定数と，$m$ に依存する減衰率 $\Gamma_m$ を，それぞれ $D,R,m$ で表せ。
\item \textbf{ゼロモード＝1 次元拡散。}$m=0$ のモード係数 $P_0(x,t)$ が $y$ で積分した周辺分布
$\displaystyle\int_0^{2\pi R}P(x,y,t)\,dy$ に等しいことを確かめ，それが満たす方程式を書け。
点源初期条件のもとで $P_0(x,t)$ を求めよ。
\item \textbf{Kaluza--Klein 塔と寿命。}$m\neq0$ のモードを点源初期条件のもとで解き，$P_m(x,t)$ を
$\Gamma_m$ と 1 次元熱核で表せ。各モードの減衰の時定数 $\tau_m$ を $R,D,m$ で求め，モードを寿命の長い順に
並べよ。最長寿命の非ゼロモードからクロスオーバー時刻 $t^\ast$ を定義し，$R,D$ で表せ。
\item \textbf{全解と 2 つの表示。}点源に対する全確率密度 $P(x,y,t)$ を書け（$x$ 部分と $y$ 部分の積になる）。
$y$ 部分（半径 $R$ の円上の熱核）を，\textbf{(i)} 運動量モード $m$ についての和，および
\textbf{(ii)} 巻きつき数（像）$w\in\mathbb{Z}$ についての和（$y$ の各像 $y-2\pi Rw$ からの自由熱核の重ね合わせ）
の\emph{2 通り}で表せ。両者は Poisson 和で結ばれる。それぞれ $t$ が小さいとき・大きいときの
どちらで速く収束するかを述べよ。
\item \textbf{低エネルギー（長時間）極限。}$t\gg t^\ast$ のとき，どのモードが生き残るかを述べ，
$P(x,y,t)$ の主要項を求めよ。これが $y$ に一様で $x$ 方向の 1 次元拡散になる（余剰次元が見えない）ことを示せ。
さらに最初の補正（最低次の非ゼロモード）の大きさの $t$ 依存性を求めよ。
\item \textbf{高エネルギー（短時間）極限。}$t\ll t^\ast$ のとき，$y$ 方向の広がり $\sqrt{2Dt}$ と円周 $2\pi R$ を
比べよ。(4)(ii) の像和で支配的になる項を述べ，$P(x,y,t)$ の主要項を求めよ。これが自由 2 次元拡散
（コンパクト性が見えない）になることを示せ。また (4)(i) のモード和で実質的に寄与する $m$ の個数を
$R,D,t$ で見積もれ。
\item \textbf{スケーリング次元と有効次元。}拡散方程式のスケール不変性から，解の $t$ 冪が
空間次元を``測る''ことを見る。
\begin{enumerate}
\item $d$ 次元拡散方程式 $\partial_t P=D\nabla^2P$ を不変に保つスケール変換
$x\to bx,\ t\to b^z t$ の指数 $z$ を求めよ。このとき $D$ はどう変換するか。
\item 規格化 $\int P\,d^dx=1$ を保つには $P$ がどう変換しなければならないか。これと (a) から，
点源解の自己相似形 $P(x,t)=t^{-\alpha}\Phi\!\bigl(x/\sqrt{4Dt}\bigr)$ の冪 $\alpha$ を空間次元 $d$ の
関数で表し，$d=1,2$ で評価せよ。
\item 本問（$\mathbb{R}\times S^1$）の原点回帰確率 $P(0,0,t)$ の $t$ 依存の冪を，$t\gg t^\ast$ と
$t\ll t^\ast$ のそれぞれで求めよ。各領域でこの冪が``見かけ上''何次元の拡散の $\alpha$ に一致するか
（有効次元 $d_{\mathrm{eff}}$）を述べ，その領域で実質的に効いている KK モードの個数と対応づけよ。
\item (c) を具体的に示せ：長時間では $K(0,t)=\frac{1}{2\pi R}\sum_m e^{-Dm^2t/R^2}$ のどの項が残り
$P(0,0,t)$ の冪が何になるか；短時間ではこの和を積分で近似して $K(0,t)$ の $t$ 冪を求め，$P(0,0,t)$ の
冪が何になるかを示せ。有効次元が UV（短時間）の $2$ から IR（長時間）の $1$ へ\emph{流れる}ことを，
効くモード数の言葉でまとめよ。
\end{enumerate}
\item \textbf{考察：運動量--巻きつき双対と次元簡約。}
\begin{enumerate}
\item (4) の 2 表示（運動量モード和 $m$ と巻きつき和 $w$）が，それぞれ長時間・短時間で支配的になる
理由を，(5)(6) の結果と結びつけて述べよ。
\item 本問の構造を Kaluza--Klein 次元簡約（コンパクト次元 $\to$ ``質量'' $\propto m/R$ のモードの塔，
低エネルギーでは質量ゼロのモードのみが効く）と対応づけよ。ここで何が「エネルギースケール」に，
何が「コンパクト化スケール」に対応するか。減衰率 $\Gamma_m$ は KK 質量とどう対応するかを述べよ。
\item この拡散系で，運動量モード $m$ と巻きつき $w$ の役割を入れ替える「双対」が成り立つとすれば，
$R$ はどのような量に置き換わると期待されるか（次元解析の範囲で）論ぜよ。
\end{enumerate}
\end{enumerate}
\end{prob}

\bigskip
{\sloppy
\noindent\textbf{解答について：}各問の完全な解答（求めるべき表式・数値を含む）は別ファイル
\texttt{exercise\_diffusion\_extra\_dimension\_solutions.tex} にある。
\par}

\end{document}
